\section*{第 \ref{ch:b1:active-margins} 章——活动大陆边缘}

\begin{solution}{exo:b1:active-margins:1}
(a) $2\pi G\,\Delta\rho\,h = 2\pi\times\num{6.674e-11}\times1700\times3000 =
\qty{2.14e-3}{m.s^{-2}}$：重力将降低约
\qty{214}{\milli\gal}。(b) 对于 \qty{3.5}{km}，
$2\pi\times\num{6.674e-11}\times1700\times3500 = \qty{2.50e-3}{m.s^{-2}}$,
\qty{250}{\milli\gal}：与观测到的低值相当，\omterm{def:b1:active-margins:active-margin}{海沟}因此远离
均衡状态。(c) 板片公式假设缺失的岩石
向四周无限延伸；宽度仅数十千米的\omterm{def:b1:active-margins:active-margin}{海沟}
在测量点周围移去的质量少得多，因此仅海水
产生的效应便小于 \qty{250}{\milli\gal}。其余亏损位于更深处：
整个板块在\omterm{def:b1:active-margins:active-margin}{海沟}下方向下弯曲，其\omterm{def:b1:layered-earth:layers}{地壳}比
\omterm{def:b1:layered-earth:layers}{地幔}轻，在其底部取代了\omterm{def:b1:layered-earth:layers}{地幔}的位置。
\end{solution}

\begin{solution}{exo:b1:active-margins:2}
(a) $\qty{82}{mm.\text{年}^{-1}} = \qty{82}{km.Ma^{-1}}$，故
\qty{10}{Ma} 内移动 \qty{820}{km}。(b) $820\sin\ang{45} = \qty{580}{km}$。(c) 到达
\qty{600}{km} 处的路径长为 $600/\sin\ang{45} = \qty{849}{km}$，
用时 $849/82 = \qty{10.3}{Ma}$。
\end{solution}

\begin{solution}{exo:b1:active-margins:3}
$g = 110/\tan\delta$：\ang{30} 时为 \qty{191}{km}，\ang{45} 时为 \qty{110}{km}，
\ang{60} 时为 \qty{64}{km}。缓倾板片上方的\omterm{def:b1:active-margins:volcanic-arc}{火山弧}
距其\omterm{def:b1:active-margins:active-margin}{海沟}最远。
\end{solution}

\begin{solution}{exo:b1:active-margins:4}
(a) 汇聚速率：$\qty{2.7e6}{m^2.\text{年}^{-1}}/\qty{5.1e7}{m} =
\qty{0.053}{m.\text{年}^{-1}}$，约 \qty{53}{mm.\text{年}^{-1}}。扩张速率：
$\qty{3.0e6}{m^2.\text{年}^{-1}}/\qty{6.7e7}{m} = \qty{0.045}{m.\text{年}^{-1}}$，约
\qty{45}{mm.\text{年}^{-1}}。在总长度较短的情况下，俯冲边界
消耗洋底的速度快于洋脊生成洋底的速度，但总消耗量较少。
(b) 其余 \qty{0.3}{km^2} 左右沿模型中其他的汇聚
边界被吸收，例如大陆碰撞和海洋内的汇聚带，
这些地方板块被缩短而非俯冲；并且
洋脊并非唯一的离散型边界：\omterm{def:b1:rifts-ridges-margins:rift}{大陆裂谷}和部分
\omterm{def:b1:rifts-ridges-margins:transform}{转换断层}也略有张开。
\end{solution}

\begin{solution}{exo:b1:active-margins:5}
(a) $\bar x = 350$，$\bar z = 290$；$x$ 的偏差依次为 $-150$、$-50$、
$50$、$150$ 以及 $z$ 的偏差为 $-140$、$-40$、$40$、$140$，因此
$S_{xx} = 2(150^2 + 50^2) = \num{50000}$，而
$S_{xz} = 2(150\times140 + 50\times40) = \num{46000}$。于是 $s = 0.92$，
$a = 290 - 0.92\times350 = \qty{-32}{km}$，\omterm{def:b1:geological-maps:attitude}{倾角}为
$\arctan 0.92 = \ang{42.6}$。(b) 当 $z = 0$ 时 $x = 32/0.92 = \qty{34.8}{km}$
距\omterm{def:b1:active-margins:active-margin}{海沟}。(c) 由端点求得斜率为 $280/300 = 0.933$，\omterm{def:b1:geological-maps:attitude}{倾角}
为 \ang{43.0}，较接近，但仅用了四次地震中的两次。
\end{solution}

\begin{solution}{exo:b1:active-margins:6}
(a) $50/\tan\ang{25} + 60/\tan\ang{45} = 107.2 + 60.0 = \qty{167.2}{km}$.
(b) 若板片全程\omterm{def:b1:geological-maps:attitude}{倾角}为 \ang{45}，则得 \qty{110}{km}：缓倾的浅部
段使\omterm{def:b1:active-margins:volcanic-arc}{弧沟间隙}增大 \qty{57}{km}。(c) 深部段为
$z = 50 + (x - 107.2)$；当 $z = 0$ 时 $x = \qty{57.2}{km}$，位于
\omterm{def:b1:active-margins:active-margin}{海沟}陆侧，与千岛线在 \qty{53}{km} 处的情况相同：缓倾的浅部段
可解释千岛线的错位。(d) $107.2 + 60/\tan\ang{60} = 107.2 + 34.6 =
\qty{141.9}{km}$.
\end{solution}

\begin{solution}{exo:b1:active-margins:7}
(a) 半径为\omterm{def:b1:geological-maps:attitude}{倾角}的一半，且 $r = R\,\theta$，其中 $\theta$ 以弧度为单位：
\ang{22.5} 与 $6371\times0.3927 = \qty{2502}{km}$（对于 \ang{45}）；\ang{15} 与
\qty{1668}{km}（对于 \ang{30}）；\ang{35} 与 \qty{3892}{km}（对于 \ang{70}）。
(b) 缓倾板片上方的\omterm{def:b1:active-margins:volcanic-arc}{火山弧}。(c) 预测值为
$44.9/2 = \ang{22.45}$；拟合值 \ang{21.9} 比
预测值小 $0.55/22.45 = \qty{2.4}{\percent}$。
\end{solution}

\begin{solution}{exo:b1:active-margins:8}
(a) $q/k$：\omterm{def:b1:active-margins:active-margin}{弧前}下方为 $0.035/2.5 = \qty{0.014}{K.m^{-1}} = \qty{14}{K.km^{-1}}$，
弧后下方为 $0.080/2.5 = \qty{32}{K.km^{-1}}$。
(b) $10 + 14\times30 = \qty{430}{\celsius}$ 和
$10 + 32\times30 = \qty{970}{\celsius}$。(c) 地表\omterm{def:b1:heat-budget:heat-flow}{热流}的一部分
由\omterm{def:b1:layered-earth:layers}{地壳}自身的放射性产生，因此穿过
更深层位的热量减少，温度梯度随深度减小：对地表梯度作线性
外推会高估深部的温度。
\end{solution}

\begin{solution}{exo:b1:active-margins:9}
(a) \qty{32}{mm.\text{年}^{-1}} 即 \qty{32}{km.Ma^{-1}}；半速率为
\qty{16}{mm.\text{年}^{-1}}。(b) $100/32 = \qty{3.1}{Ma}$。(c) 速率未必
一直保持不变；宽度的一部分
可能来自扩张开始前弧的伸展与裂解；
且洋脊轴可能发生跳跃或不对称扩张。
\end{solution}

\begin{solution}{exo:b1:active-margins:10}
(a) 路径长为 $600/\sin\ang{45} = \qty{849}{km}$，以
\qty{82}{km.Ma^{-1}} 的速度行进，用时 \qty{10.3}{Ma}，即 \qty{3.27e14}{s}。
(b) $\sqrt{\num{e-6}\times\num{3.27e14}} = \qty{1.8e4}{m}$，约
\qty{18}{km}. (c) $(\qty{5e4}{m})^2/\qty{e-6}{m^2.s^{-1}} = \qty{2.5e15}{s}$,
约 \qty{79}{Ma}。(d) 热量仅到达板片两侧各
外缘的 \qty{18}{km}：其核心抵达 \qty{600}{km} 处时，几乎与离开
\omterm{def:b1:active-margins:active-margin}{海沟}时一样冷，比周围\omterm{def:b1:layered-earth:layers}{地幔}更冷更强，这正是它
仍能在该深度发生地震的原因。
\end{solution}

\begin{solution}{exo:b1:active-margins:11}
(a) 在坡脚后方距离 $L$ 处，地表高出 $L\tan\alpha$
坡脚水平面，底面位于其下方 $L\tan\beta$：\omterm{def:b1:active-margins:accretionary}{增生楔}后缘的
高为 $L(\tan\alpha + \tan\beta)$，三角形面积为
$\frac12 L\times L(\tan\alpha + \tan\beta)$. (b) $\tan\ang{4} + \tan\ang{6} =
0.175$。\omterm{def:b1:active-margins:accretionary}{增生楔}接收的剖面：$\qty{1}{km}\times\qty{40}{km.Ma^{-1}} =
\qty{40}{km^2.Ma^{-1}}$（面积），故 $A = 40t$，且
$L = \sqrt{2A/0.175}$：\qty{1}{Ma} 后为 \qty{21.4}{km}，\qty{4}{Ma} 后为
\qty{42.8}{km}，\qty{10}{Ma} 后为 \qty{67.6}{km}。$L$ 随
$\sqrt t$ 增长：时间为四倍则宽度为两倍。(c) $A$ 减半，$L$
除以 $\sqrt2$：\qty{15.1}{km}、\qty{30.2}{km} 和 \qty{47.8}{km}。
\end{solution}

\begin{solution}{exo:b1:active-margins:12}
(a) 由 $\partial F/\partial a = -2\sum(z_i - a - s x_i) = 0$ 得
$n\bar z = na + sn\bar x$，即 $\bar z = a + s\bar x$。(b) 每个 $z_i$ 与
$\bar z$ 同时增加 \qty{10}{km}，故偏差 $z_i - \bar z$ 与 $S_{xz}$
保持不变：斜率不变，而 $a = \bar z - s\bar x$ 增加
\qty{10}{km}。(c) 偏差 $x_i - \bar x$ 不变，故斜率
不变；$\bar x$ 减少 20，故 $a$ 增加 $20s$。(d) 均不改变：
\omterm{def:b1:geological-maps:attitude}{倾角}只取决于斜率。所有数据点共有的误差只会使
直线平移，而不会使其转动。
\end{solution}

\begin{solution}{pb:b1:active-margins:1}
\textbf{1.} 八个点近似位于一条\omterm{def:b1:geological-maps:attitude}{倾角}约为
\ang{45} 的直线上，自距\omterm{def:b1:active-margins:active-margin}{海沟} \qty{155}{km}、深 \qty{96}{km} 处起，
至距\omterm{def:b1:active-margins:active-margin}{海沟} \qty{609}{km}、深 \qty{573}{km} 处止。

\textbf{2.} $\sum x_i = 3152$，$\sum z_i = 2751$，$\sum x_i^2 = \num{1442910}$，
$\sum x_i z_i = \num{1290548}$；$\bar x = \qty{394.0}{km}$，
$\bar z = \qty{343.875}{km}$。

\textbf{3.} $S_{xx} = 1442910 - 8\times394^2 = \num{201022}$，且
$S_{xz} = 1290548 - 8\times394\times343.875 = \num{206654}$。
于是 $s = 206654/201022 = 1.028$，且 $a = 343.875 - 1.02802\times394 =
\qty{-61.2}{km}$: $z = -61.2 + 1.028\,x$.

\textbf{4.} $\arctan1.028 = \ang{45.8}$.

\textbf{5.} 在\omterm{def:b1:active-margins:active-margin}{海沟}向陆一侧 $x = 61.2/1.028 = \qty{59.5}{km}$ 处，
以此\omterm{def:b1:geological-maps:attitude}{倾角}自\omterm{def:b1:active-margins:active-margin}{海沟}下插的板片应在该处抵达地表，
但它在\omterm{def:b1:active-margins:active-margin}{海沟}处即已抵达，故近\omterm{def:b1:active-margins:active-margin}{海沟}处\omterm{def:b1:geological-maps:attitude}{倾角}较缓，深处则变陡。

\textbf{6.} 残差为 $-2.1$，$2.0$，$-7.3$，$5.3$，$7.1$，$2.3$，
$-15.0$ 与 $8.1$ \unit{km}。事件 7 位于直线上方 \qty{15}{km}，事件 8
位于其下方 \qty{8}{km}；其余各点均在距其约 \qty{7}{km} 以内，
仅为板片厚度的一小部分。

\textbf{7.} $\ang{320.9} - \ang{300} = \ang{20.9}$。垂直于
\omterm{def:b1:active-margins:active-margin}{海沟}方向：$87.4\cos\ang{20.9} = \qty{81.6}{mm.\text{年}^{-1}}$；平行方向：
$87.4\sin\ang{20.9} = \qty{31.2}{mm.\text{年}^{-1}}$。汇聚为斜向。

\textbf{8.} \qty{81.6}{mm.\text{年}^{-1}} 即 \qty{81.6}{km.Ma^{-1}}，故 \qty{816}{km}
\qty{10}{Ma}.

\textbf{9.} 路径为 $632/\sin\ang{45} = \qty{894}{km}$，历时
$894/81.6 = \qty{11.0}{Ma}$.

\textbf{10.} $122 + 11 = \qty{133}{Ma}$.

\textbf{11.} $t = \qty{10.95}{Ma} = \qty{3.46e14}{s}$，且
$\sqrt{\num{e-6}\times\num{3.46e14}} = \qty{1.86e4}{m}$，约 \qty{19}{km}
距每个表面：远小于 \qty{100}{km} 厚板片中央的 \qty{50}{km}，
其核心仍很冷。

\textbf{12.} $q/k$：老洋底之下为 \qty{17.2}{K.km^{-1}}，
\omterm{def:b1:active-margins:active-margin}{弧前}之下为 \qty{14.0}{K.km^{-1}}，弧后之下为
\qty{32.0}{K.km^{-1}}。

\textbf{13.} 在\omterm{def:b1:active-margins:active-margin}{弧前}之下，冷板片仅位于地下数十千米深处；
它从上方岩石中吸收热量，因此到达
地表的热量减少。

\textbf{14.} $h = \qty{2.60e-3}{m.s^{-2}}/(2\pi\times\num{6.674e-11}\times1700)
= \qty{3647}{m}$，约 \qty{3.6}{km}，与
\omterm{def:b1:active-margins:active-margin}{海沟}低于周围海底的整个深度相当，甚至略大（该深度约为 \qty{3.5}{km}）。
\omterm{def:b1:active-margins:active-margin}{海沟}缺失的质量未得到任何补偿。

\textbf{15.} 密度差为 $3300 - 1000 = \qty{2300}{kg.m^{-3}}$，
且 $2\pi\times\num{6.674e-11}\times2300\times3500 = \qty{3.38e-3}{m.s^{-2}}$，
\qty{338}{\milli\gal}。板片公式将缺失质量计入任意距离；
而真实\omterm{def:b1:active-margins:active-margin}{海沟}仅在宽数十千米的带状区域内下弯，
因此其亏损较小，约为 \qty{260}{\milli\gal}，介于
两种板片估算值之间。

\textbf{16.} 弯曲板块的强度，它像弹性板块那样使\omterm{def:b1:active-margins:active-margin}{海沟}下陷、
外隆起上拱，以及致密板片
下沉进入\omterm{def:b1:layered-earth:layers}{地幔}时的拉力。

\textbf{17.} $-61.2 + 1.028\times164 = \qty{107.4}{km}$.

\textbf{18.} 是：\qty{107}{km} 介于 72 与 \qty{173}{km} 之间，
仅比平均值 \qty{105}{km} 深约 \qty{2}{km}。

\textbf{19.} $x = (100 + 61.2)/1.028 = \qty{156.8}{km}$.

\textbf{20.} $\ang{45.8}/2 = \ang{22.9}$，且
$6371\times22.9\times\pi/180 = \qty{2546}{km}$.

\textbf{21.} 拟合值 \ang{21.9}，小 \ang{1.0}，约
\qty{4}{\percent}：弯曲壳层能很好地描述千岛弧的
曲率。

\textbf{22.} $x = (105 + 61.2)/1.028 = 166.2/1.028 = \qty{161.7}{km}$.

\textbf{23.} $164 - 161.7 = \qty{2.3}{km}$，比观测到的
火山前缘更靠近\omterm{def:b1:active-margins:active-margin}{海沟}；$2.3/164 = \qty{1.4}{\percent}$：千岛火山
位于比平均深度略深的板片之上，正如第 18 题所示。

\textbf{24.} 取整到千米，千岛板片直线预测的
\omterm{def:b1:active-margins:volcanic-arc}{弧沟间隙}为 \textbf{\qty[text-series-to-math=true]{162}{km}}，适用于
板片深 \qty{105}{km} 处之上的火山（观测值：\qty{164}{km}）。
\end{solution}
